% Research fragment: independently audited 2026-10-08.
% The parent supplies theorem environments, amsmath/amssymb, and sources.bib.

\section{An exponential full-complex checkpoint inside explicit unknots}
\label{sec:unknot-checkpoint-lower-bound}

The Boolean connection-rank obstruction concerns a representation that
answers every admitted continuation.  The next obstruction is different:
it applies to the explicit size of a complex that retains the complete
Bar--Natan homotopy type of one specified prefix.  It holds even when the
input is promised to be an unknot.

\subsection{A reduced closure functor needs no exactness hypothesis}

Fix a field $\mathbb F$ and the ordinary Khovanov Frobenius algebra
$A=\mathbb F[X]/(X^2)$.  Work in the additive Bar--Natan category with
fixed boundary, using the standard specialization if a universal
deformation is present.  We count actual summands, with multiplicity,
of circle-free crossingless matchings in a bounded, explicitly delooped
complex.  A quantum or homological shift does not change this count.

\begin{lemma}[Closure lower bound]
\label{lem:reduced-closure-size}
Let $P$ be a tangle with $2r>0$ boundary points.  Fix a crossingless
closure $\Gamma$ for which $\Gamma(P)$ is a knot.  Suppose that $C$ is
homotopy equivalent to the full Bar--Natan complex of $P$ and that all
objects appearing in $C$ are shifts of circle-free crossingless
matchings.  If $M(C)$ is the total number of summands over all chain
degrees, counted with multiplicity, then
\[
 M(C)\ \ge\ 2^{1-r}\,
 \dim_{\mathbb F}\widetilde{Kh}(\Gamma(P);\mathbb F)
 \ \ge\ 2^{1-r}\det(\Gamma(P)).
\]
In particular, for a four-ended tangle and numerator closure,
$M(C)\ge \det(N(P))/2$.
\end{lemma}

\begin{proof}
Choose a point $*$ in the interior of one of the external closure arcs.
Closing a cobordism adds the fixed product $\Gamma\times[0,1]$.
Consequently the track $*\times[0,1]$ lies on the resulting cobordism
and joins its distinguished input and output circles.  There is no
requirement that the original tangle category carry a basepoint.

Let $F$ be the ordinary Khovanov TQFT.  Give $F(\Gamma(S))$ its
$A$-module structure by multiplication on the circle containing $*$.
Every map obtained by closing a tangle cobordism is $A$-linear: a
multiplication by $X$ can be moved along the fixed marked track from the
incoming end to the outgoing end.  Algebraically this follows from
commutativity of multiplication and the Frobenius identity
\[
 \Delta(Xa)=(X\otimes1)\Delta(a)
             =(1\otimes X)\Delta(a).
\]
Births and deaths away from the marked track commute with this action
as well.  These identities apply to composite cobordisms and their
linear combinations.  The usual sphere, torus, and neck-cutting
relations are already respected by $F$.

Thus the assignment
\[
 R_\Gamma(S)=F(\Gamma(S))\otimes_A A/(X)
\]
and the induced maps define an additive functor on the original,
unpointed tangle category.  Equivalently, up to a grading shift, one
can restrict to the image of multiplication by $X$ on the marked
circle.  All modules before reduction are free as $A$-modules, so
the quotient and image descriptions are naturally identified by
\mbox{$[a]\mapsto Xa$}.

An additive functor preserves the equations defining a chain homotopy.
For example, $gf-1=dh+hd$ implies
$R_\Gamma(g)R_\Gamma(f)-1
=R_\Gamma(d)R_\Gamma(h)+R_\Gamma(h)R_\Gamma(d)$.
It therefore preserves chain homotopy equivalences.  No exactness of
$-\otimes_A A/(X)$ is asserted or needed.  Applied to the cube of
resolutions of $P$, the functor gives the reduced Khovanov complex of
$\Gamma(P)$, up to harmless grading shifts.  This is compatible with
the planar-composition and TQFT constructions in
\cite[Theorem~2 and Section~7]{BarNatanTangles} and with the reduced
complex in \cite[Section~3]{KhovanovPatterns}.

A circle-free matching on $2r$ points has $r$ arcs.  Gluing it to the
$r$ closure arcs produces $c$ circles with $1\le c\le r$.
Hence
\[
 \dim_{\mathbb F}R_\Gamma(S)=2^{c-1}\le2^{r-1}.
\]
It follows that the total vector-space dimension of
$R_\Gamma(C)$ is at most $2^{r-1}M(C)$.  Total homology dimension is
at most total chain dimension, and homotopy equivalence identifies
its homology with $\widetilde{Kh}(\Gamma(P);\mathbb F)$.

Finally, the graded Euler characteristic of reduced Khovanov homology
is the normalized Jones polynomial.  Evaluating its grading variable
at $i$ corresponds to evaluating the conventional Jones variable at
$-1$; all grading monomials have absolute value one.  The triangle
inequality gives
\[
 \dim_{\mathbb F}\widetilde{Kh}(L;\mathbb F)
 \ge |V_L(-1)|=|\Delta_L(-1)|=\det(L).
\]
This is the elementary reduced-homology inequality recorded in
\cite[Section~3]{KhovanovPatterns}.  The same Euler-characteristic
argument works over every field, since the integer graded Euler
characteristic of the finite complex is independent of the field.
Combining the inequalities proves the result.
\end{proof}

\begin{remark}[An independent unreduced proof over $\mathbb F_2$]
The four-ended bound also follows without constructing a reduced
closure functor.  An ordinary numerator closure sends each
circle-free matching to at most two circles and thus at most four
scalar generators.  Over $\mathbb F_2$, the unreduced Khovanov
homology is the direct sum of two grading shifts of the reduced
homology \cite[Corollary~3.2.C]{ShumakovitchTorsion}.  Therefore
\[
 4M(C)\ge\dim Kh(N(P);\mathbb F_2)
 =2\dim\widetilde{Kh}(N(P);\mathbb F_2)
 \ge2\det(N(P)).
\]
This supplies the same lower bound $M(C)\ge\det(N(P))/2$.
\end{remark}

\subsection{Pell growth at a four-ended unknot checkpoint}

Use the matrices and tangle conventions of
\eqref{eq:rational-two-twist-generators}.  For $k\ge0$, set
\[
 G_k=(AB)^k,
 \qquad \binom{p_k}{q_k}=G_k\binom11,
 \qquad \binom{r_k}{s_k}=G_k\binom32,
\]
and choose the standard rational tangle diagrams
\[
 P_k=T(p_k/q_k),\qquad Q_k=-T(r_k/s_k),
 \qquad K_k=N(P_k+Q_k).
\]

\begin{theorem}[Explicit unknots with an exponential intrinsic checkpoint]
\label{thm:unknot-exponential-checkpoint}
There are diagrams $K_k$ of the unknot with $N_k=8k+4$ displayed
crossings and a specified four-ended prefix $P_k$ such that every
explicitly delooped complex $C_k$ retaining the complete Bar--Natan
homotopy type of $P_k$ has
\[
 M(C_k)\ge\frac{p_k}{2}
 \ge\frac12(3+2\sqrt2)^k
 =\frac12(3+2\sqrt2)^{(N_k-4)/8}.
\]
Here
\[
 p_0=1,\quad p_1=7,\quad
 p_{k+2}=6p_{k+1}-p_k,
\]
and, with $\lambda=3+2\sqrt2$,
\[
 p_k=\frac{1+\sqrt2}{2}\lambda^k
       +\frac{1-\sqrt2}{2}\lambda^{-k}.
\]
In particular, the lower bound is exponential in the size of these
displayed unknot diagrams, although the rational arithmetic summaries
have $O(k)$ bits and decide their closure in $O(k^2)$ bit operations.
\end{theorem}

\begin{proof}
Both columns are primitive and
\[
 p_ks_k-q_kr_k
 =\det(G_k)\det\begin{pmatrix}1&3\\1&2\end{pmatrix}
 =-1.
\]
The rational gluing criterion
\eqref{eq:rational-unknot-gluing} therefore proves that $K_k$ is an
unknot.  The relevant classical criterion is
\cite[Theorem~5]{KauffmanLambropoulouHard}.

The prescribed construction uses $4k+1$ crossings for $P_k$ and
$4k+3$ for $Q_k$, and numerator caps introduce none.
For an explicit continued-fraction realization, $P_0=[1]$ and, for
$k\ge1$,
\[
 P_k=[\underbrace{2,\ldots,2}_{2k-1},3],
 \qquad
 -Q_k=[\underbrace{2,\ldots,2}_{2k},1,2].
\]
The sums of the absolute coefficients are exactly $4k+1$ and
$4k+3$.  These are counts in the displayed input diagrams, not
minimal crossing numbers of their knot type.

The matrix
\[
 AB=\begin{pmatrix}5&2\\2&1\end{pmatrix}
\]
has trace six and determinant one.  Its characteristic equation
gives the asserted recurrence.  Its eigenvalues are
$\lambda$ and $\lambda^{-1}$, and the initial values give the stated
closed formula.  Since
\[
 p_k-\lambda^k
 =\frac{\sqrt2-1}{2}(\lambda^k-\lambda^{-k})\ge0,
\]
we have $p_k\ge\lambda^k$.

Moreover, $G_k\equiv I\pmod2$, so $p_k$ and $q_k$ are odd.
Thus $N(P_k)$ is a rational knot and its determinant is $p_k$.
The determinant assertion is the usual primitive fraction formula
for rational numerator closure; it can also be computed from its
lens-space double branched cover, whose first homology has order
$p_k$.  The tangle fraction and numerator/denominator determinant
interpretation are reviewed in
\cite[Section~4.1, Remark~2]{KauffmanLambropoulouFractions}.
Apply Lemma~\ref{lem:reduced-closure-size} with $r=2$ to obtain
$M(C_k)\ge p_k/2$.  Since $M(C_k)$ is integral and $p_k$ is odd, the
rounded bound is $(p_k+1)/2$.

Each matrix update uses additions and doubling of $O(k)$-bit
integers.  The two slope columns and their determinant consequently
require $O(k)$ storage bits and $O(k^2)$ schoolbook bit operations,
as in Theorem~\ref{thm:rational-arithmetic-representation}.
\end{proof}

\begin{corollary}[Stronger bound when the quantum grading is retained]
\label{cor:graded-checkpoint-size}
For the ordinary graded Bar--Natan complex, the conclusion of
Theorem~\ref{thm:unknot-exponential-checkpoint} strengthens to
\[
 M(C_k)\ge p_k+q_k,
 \qquad q_k=\frac{\lambda^k+\lambda^{-k}}2.
\]
Indeed the two individual circle-free matching types must occur at
least $p_k$ and $q_k$ times, respectively, counting all shifts and
chain degrees.  More generally, for any number of ends a graded
version of Lemma~\ref{lem:reduced-closure-size} gives
$M(C)\ge\det(\Gamma(P))$.
\end{corollary}

\begin{proof}
Use $z$ for the quantum grading variable to distinguish it from the
fraction denominator $q_k$.  A matching whose closure has $c$
circles contributes, up to a grading monomial,
\[
 (z+z^{-1})^{c-1}
\]
to the graded dimension of the reduced closure.  At $z=i$ this is
zero if $c>1$ and has absolute value one if $c=1$.  The reduced
graded Euler characteristic therefore yields
\[
 \det(\Gamma(P))\le
 \#\{\hbox{summands whose $\Gamma$-closure is one circle}\}.
\]
This proves the general assertion.  It only uses graded homotopy
equivalence and evaluation of the graded Euler characteristic; it
does not assume thinness of Khovanov homology.

For four ends there are exactly two circle-free matchings, the
horizontal and vertical ones.  Numerator closure gives two circles
on the horizontal matching and one on the vertical matching;
denominator closure interchanges these roles.  Let $M_H$ and $M_V$
be their total multiplicities.  The numerator estimate gives
$M_V\ge p_k$.  The denominator closure of $P_k$ is also a rational
knot because $q_k$ is odd, and has determinant $q_k$.  Applying the
same estimate to denominator closure gives $M_H\ge q_k$.  Thus
$M(C_k)=M_H+M_V\ge p_k+q_k$.

The recurrence for $q_k$ is the same trace-six recurrence as for
$p_k$, with initial values $q_0=1$ and $q_1=3$, giving the displayed
formula.  One orientation of the prefix need not extend to both test closures.
Use the orientation-free bracket complex first; changing strand orientations
changes the normalized complex only by global homological and quantum shifts.
Their evaluations have absolute value one. Thus both estimates apply to
the same matching multiplicities, despite independent closure orientations.
\end{proof}

\paragraph{Scope of the obstruction.}
This theorem applies at the specified prefix checkpoint and to
representations preserving its full homotopy type in the stated
explicit basis.  It is not a lower bound for every scan order or
every unknot-recognition algorithm.  An algorithm can use information
about the actual continuation, absorb parts of it earlier, or retain
the rational slope in a compact arithmetic form.  The conclusion
also does not charge an exponentially repeated object encoded by a
binary multiplicity as exponentially many stored words: it rules
out a small \emph{expanded} full complex, and any use of compressed
multiplicities or other compressed objects must prove that all
required operations on that representation remain efficient.

The distinction from the Boolean identity-matrix theorem is precise.
That theorem rules out a small exact linear representation of all
Boolean continuation answers on its family.  The present theorem
rules out a small expanded full homotopy representative for one
sequence of prefixes inside unknots.  Both leave room for the
rank-two signed arithmetic pairing followed by a nonlinear test.
